\documentclass{article}

\usepackage[utf8]{inputenc}
\usepackage[T1]{fontenc}
\usepackage{amsmath,amssymb,amsthm}
\usepackage{mathtools}
\usepackage{hyperref}
\usepackage{booktabs}

\makeatletter
\newcommand{\citep}[2][]{%
  \begingroup
  \def\@tmpkey{#2}%
  \ifx\@tmpkey\@empty\else
    (\ifx\@empty#1\@empty\else #1, \fi
     \@nameuse{bibkey@\@tmpkey})%
  \fi
  \endgroup
}
\expandafter\def\csname bibkey@lasser2026micro\endcsname{Microfoundation}
\expandafter\def\csname bibkey@lasser2026exo\endcsname{Exogenous~Verification}
\expandafter\def\csname bibkey@lasser2026phase\endcsname{Phase~Redundancy}
\makeatother

\newtheorem{definition}{Definition}
\newtheorem{lemma}{Lemma}
\newtheorem{theorem}{Theorem}
\newtheorem{remark}{Remark}
\newtheorem{assumption}{Assumption}
\newtheorem{proposition}{Proposition}

\newcommand{\Bclip}{B_{\mathrm{clip}}}
\newcommand{\Kch}{K_{\mathrm{ch}}}
\newcommand{\Aadv}{\mathcal{A}_{\mathrm{adv}}}
\newcommand{\Pproxy}{P^{\mathrm{proxy}}}
\newcommand{\Ttruth}{T^{\mathrm{truth}}}
\newcommand{\epsnonresCoev}{\varepsilon_{\mathrm{coev}}^{\mathrm{nonres}}}
\newcommand{\Nev}{N_{\mathrm{ev}}}
\newcommand{\rhogap}{\rho_{\mathrm{gap}}}
\newcommand{\lammax}{\lambda_{\max}}
\newcommand{\taumeta}{\tau_{\mathrm{meta}}}
\newcommand{\Kchmulti}{K_{\mathrm{ch}}^{\mathrm{multi}}}

\title{C7 closed-form derivation, v2: \\
Lipschitz cross-channel sensitivity and cumulative cascade-window bound}
\author{Paper 10 working draft for codex review}
\date{}

\begin{document}

\maketitle

\section{Goal and changes from v1}

\subsection{Goal}

Derive (C7) bounded co-evolution from primitive operational
parameters, so that $\bar{M}$ is no longer a free-floating
deployment-class assumption but a corollary of (C5.HOEFF),
(C5.MULT), and a structural property of the verification
protocol's channel partition.

\subsection{What v1 got wrong (per codex review)}

\begin{enumerate}
\item \textbf{Definition inconsistency (P0).} v1 defined
  $M_{j \to j'}$ as total drift of $\mu_{j'}$ from baseline, but
  Step~4 of the proof excluded $j'$-only drift from coupling.  Under
  v1's definition, an adversary perturbing only $j'$-only actions
  would set $|\mu_{j'}(q) - \mu_{j'}(q_0)| > 0$ even when channels
  are action-disjoint, contradicting the closed-form claim
  $\bar{M} = 0$ in that case.

\item \textbf{Per-step vs.\ cumulative ambiguity (P0).} v1 proved a
  per-step bound but the operational target (Audit~4 ``within
  $\taumeta$'' / Microfoundation $K_{\mathrm{coev}}$ aggregation)
  is cumulative.  v1 left this gap unresolved.

\item \textbf{Step~4 false equality (P1).} v1's claim that ``$q$ and
  $q'$ agreeing on $A \setminus A_{j,j'}$ implies $\mu_j(q) =
  \mu_j(q')$'' is false; the agreement region needs to exclude $A_j$
  as a whole, not just $A_{j,j'}$.

\item \textbf{LLR support vs.\ structural incidence (P1).} v1
  identified ``action $a$ affects channel $j$'' with $\ell^{(j)}(a)
  \neq 0$, which is not generally valid: a channel's LLR can be
  nonzero on outcomes that are not operationally causal for that
  channel.

\item \textbf{$Q^*$ defined only over admissible adversaries
  (P1).} v1's Step~3 used $q_0(A_{j,j'}) \leq Q^*$, but $Q^*$ was
  defined only over $\mathcal{A}_{\mathrm{adv}}$, not the baseline.

\item \textbf{$Q^*$ intensivity not implied by (C5.MULT) (P1).}
  v1 claimed $Q^*$ is intensive in $|P|$ from (C5.MULT) alone, but
  (C5.MULT) fixes channel cardinality without preventing
  adversaries from concentrating mass on shared actions as $|P|$
  grows.

\item \textbf{Single $q_0$ compatibility (P1).} v1 assumed a single
  baseline $q_0$ corresponds to every per-channel $p_0^{(j)}$
  without stating the compatibility condition.

\item \textbf{Pairwise terminology (P2).} v1 prose said ``more than
  one channel'' but the math used pairwise overlap.
\end{enumerate}

\subsection{Structural decisions for v2}

\begin{itemize}
\item Replace v1's total-drift definition with an
  \textbf{L1-normalized directional sensitivity} (a Lipschitz
  formulation that avoids the denominator-vanishing problem).
\item Move to \textbf{cumulative} bound directly, with explicit
  $\lammax \cdot \taumeta$ event-count factor; this matches Audit~4's
  ``within $\taumeta$'' phrasing.
\item Introduce a \textbf{structural channel-incidence map} $I_j$
  separate from LLR support.
\item Define $Q^*$ over the union $\{q_0\} \cup \mathcal{A}_{\mathrm{adv}}
  \cup \mathcal{A}_{\mathrm{adv}}^{\mathrm{env}}$; use the
  difference quantity $Q_\Delta^*$ for the bound.
\item Add an explicit \textbf{overlap-mass stability sub-clause}
  (C5.OVL) requiring $\sup_{|P|} \sup_q q(A_{j,j'}) \leq Q_{\max} <
  1$.
\item Add a \textbf{baseline compatibility} statement: $q_0$'s
  channel projections equal $p_0^{(j)}$.
\end{itemize}

\section{Setup}

\subsection{Verification ledger and action space}

The deployment's verification ledger records events from a
finite-or-countable action space $A$ (event types under the
ledger's classification policy).  An \emph{adversary} is identified
with a probability measure $q$ on $A$ (the per-step probability of
each event type under that adversary's strategy).  Let $\mathcal{Q}
:= \{q_0\} \cup \mathcal{A}_{\mathrm{adv}} \cup
\mathcal{A}_{\mathrm{adv}}^{\mathrm{env}}$ denote the set of all
admissible action distributions, including the baseline $q_0$.

\subsection{Channel structure}

The deployment monitors $\Kch$ channels indexed by $j \in \{1,
\ldots, \Kch\}$.  Each channel has:
\begin{itemize}
\item A \textbf{structural incidence map} $I_j: A \to \{0, 1\}$
  determined by the verification protocol's classification policy:
  $I_j(a) = 1$ iff event $a$ is operationally classified into
  channel $j$.  The set $A_j := \{a \in A : I_j(a) = 1\}$ is
  channel $j$'s incidence set, fixed before deployment under
  (C5.MULT).
\item A baseline distribution $p_0^{(j)}$ over the events
  classified into $A_j$.
\item A least-favorable adversarial alternative $p_1^{(j)}$ at
  threshold shift $\eta_j$.
\item A clipped per-step log-likelihood-ratio function
  $\ell^{(j)}: A \to [-\Bclip, \Bclip]$, with $\ell^{(j)}(a) = 0$
  for $a \notin A_j$ by structural extension (zero-outside-incidence
  convention).
\end{itemize}

\begin{remark}[Incidence vs.\ LLR support]
\label{rem:incidence}
The structural incidence map $I_j$ is a separate object from
$\ell^{(j)}$'s support.  $I_j$ is \emph{operationally} fixed: it
encodes the verification protocol's channel-classification policy,
and identifies which events ``count toward'' channel $j$'s tests.
$\ell^{(j)}$ is the formal LLR function used in SPRT detection.
The convention $\ell^{(j)}(a) = 0$ for $a \notin A_j$ enforces that
operational classification respects formal channel structure: events
outside $A_j$ contribute nothing to channel $j$'s SPRT statistic.
The opposite convention (where $\ell^{(j)}$ may be nonzero off
$A_j$) corresponds to verification-protocol designs where channel
boundaries are LLR-derived rather than classification-derived; we do
not consider such designs here.
\end{remark}

\subsection{Pairwise shared-action sets}

For ordered pairs $(j, j')$ with $j \neq j'$, define:
\begin{align}
  A_{j, j'} &:= A_j \cap A_{j'} \quad \text{(actions classified into
    both channels)}.
\end{align}
$A_{j, j'} = \emptyset$ when channels $j$ and $j'$ are
\textbf{action-disjoint} by classification.

\subsection{Per-step expected drift}

For an adversary $q \in \mathcal{Q}$, the per-step expected drift
in channel $j$ is
\begin{equation}
\label{eq:drift-def}
  \mu_j(q) \;:=\; \mathbb{E}_{a \sim q}[\ell^{(j)}(a)]
  \;=\; \sum_{a \in A_j} q(a) \cdot \ell^{(j)}(a),
\end{equation}
the second equality by the zero-outside-incidence convention.

\subsection{Baseline compatibility}

\begin{assumption}[Baseline ledger compatibility]
\label{ass:baseline-compat}
There exists a baseline ledger distribution $q_0$ on $A$ such that
for every channel $j \in \{1, \ldots, \Kch\}$, the conditional
distribution $q_0(\cdot \mid A_j)$ equals the per-channel baseline
$p_0^{(j)}$, and $\mu_j(q_0) \approx 0$ (the LLR has near-zero
expectation under the per-channel baseline by SPRT design).
\end{assumption}

\section{Operational definition of \texorpdfstring{$M_{j \to j'}$}{M\_jj'}}

\subsection{j-supported perturbations}

A \emph{$j$-supported signed measure} is a function $\Delta: A \to
\mathbb{R}$ with $\Delta(a) = 0$ for $a \notin A_j$ and
$\sum_{a \in A_j} \Delta(a) = 0$.  Such $\Delta$ represents an
adversary's deviation from baseline restricted to channel $j$'s
incidence set.

The set of \textbf{admissible $j$-perturbations} is
\begin{equation}
  \mathcal{D}_j \;:=\; \{ \Delta : A \to \mathbb{R} \mid \Delta
  \text{ is $j$-supported}, \;\; q_0 + \Delta \in \mathcal{Q},
  \;\; q_0 + \Delta \geq 0 \}.
\end{equation}
For each admissible $\Delta \in \mathcal{D}_j$, the adversary
$q := q_0 + \Delta$ deviates from baseline only on $A_j$, and the
induced per-channel drift shifts are
\begin{align}
  \Delta\mu_j[\Delta] &:= \mu_j(q_0 + \Delta) - \mu_j(q_0)
    = \sum_{a \in A_j} \Delta(a) \cdot \ell^{(j)}(a), \\
  \Delta\mu_{j'}[\Delta] &:= \mu_{j'}(q_0 + \Delta) - \mu_{j'}(q_0)
    = \sum_{a \in A_{j, j'}} \Delta(a) \cdot \ell^{(j')}(a),
\end{align}
the second equality because $\Delta$ is supported on $A_j$ and
$\ell^{(j')}$ vanishes off $A_{j'}$, so only $A_j \cap A_{j'} =
A_{j,j'}$ contributes to channel $j'$'s drift.

\subsection{Per-step coupling magnitude}

\begin{definition}[Per-step coupling magnitude, L1-normalized]
\label{def:M-formal-v2}
For $j \neq j'$, the \emph{per-step coupling magnitude} of channel
$j$ on channel $j'$ is the maximum $j'$-drift response per
unit-magnitude $j$-perturbation:
\begin{equation}
\label{eq:M-def-v2}
  M_{j \to j'}^{\mathrm{step}}(\mathcal{D})
  \;:=\; \sup_{\Delta \in \mathcal{D}_j \,:\, \|\Delta\|_1 \leq 1}
  \big| \Delta\mu_{j'}[\Delta] \big|.
\end{equation}

The \emph{per-step deployment coupling magnitude} is
$\bar{M}^{\mathrm{step}}(\mathcal{D}) := \max_{j \neq j'}
M_{j \to j'}^{\mathrm{step}}(\mathcal{D})$.
\end{definition}

\begin{remark}[Why this is the Lipschitz formulation]
\label{rem:lipschitz}
Definition~\ref{def:M-formal-v2} captures the audit's
``unit-perturbation propagation'' reading directly: a
unit-$\ell_1$-magnitude $j$-perturbation produces a $j'$-drift
shift bounded by $M_{j \to j'}^{\mathrm{step}}$.  This is the
operator norm of the linear map $\Delta \mapsto \Delta\mu_{j'}$
from the space of $j$-supported signed measures (with $\ell_1$ norm)
to $\mathbb{R}$.

In differential terms: if the adversary's strategy is parametrized
by $\theta$ along a curve in $\mathcal{D}_j$, with
$\frac{\partial}{\partial \theta} (q_0 + \Delta_\theta)|_{A_j}$
having unit $\ell_1$ norm,
then $\big|\frac{\partial \mu_{j'}}{\partial \theta}\big|
\leq M_{j \to j'}^{\mathrm{step}}$.

By contrast, the v1 definition (total $j'$-drift from baseline
under any adversary) confused total drift with cross-channel
sensitivity: an adversary perturbing only $A_{j'} \setminus A_j$
would set total $j'$-drift positive but the cross-channel
sensitivity to $A_j$ perturbations is zero.  The L1-normalized
formulation isolates only $A_j$-perturbation responses.
\end{remark}

\section{Per-step closed-form bound}

\begin{lemma}[Per-step bound on $M_{j \to j'}^{\mathrm{step}}$]
\label{lem:per-step}
Under (C5.HOEFF) (per-step LLR clipped to $[-\Bclip, \Bclip]$) and
the structural-incidence convention (Section 2.2), for every $j
\neq j'$:
\begin{equation}
\label{eq:per-step-bound}
  M_{j \to j'}^{\mathrm{step}}(\mathcal{D}) \;\leq\; \Bclip,
\end{equation}
with equality possible only when $A_{j, j'}$ is non-empty and the
LLR magnitudes align favorably.  Furthermore, when the channels
are action-disjoint ($A_{j, j'} = \emptyset$),
$M_{j \to j'}^{\mathrm{step}}(\mathcal{D}) = 0$.
\end{lemma}

\begin{proof}
By Definition~\ref{def:M-formal-v2},
\[
  \big| \Delta\mu_{j'}[\Delta] \big|
  = \big| \sum_{a \in A_{j, j'}} \Delta(a) \ell^{(j')}(a) \big|
  \leq \sum_{a \in A_{j, j'}} |\Delta(a)| \cdot |\ell^{(j')}(a)|
  \leq \Bclip \sum_{a \in A_{j, j'}} |\Delta(a)|
  \leq \Bclip \cdot \|\Delta\|_1
  \leq \Bclip
\]
under $\|\Delta\|_1 \leq 1$.  Taking the supremum gives
\eqref{eq:per-step-bound}.

For the disjoint case, $A_{j, j'} = \emptyset$ implies the sum is
empty, hence $\Delta\mu_{j'}[\Delta] = 0$ for all admissible
$\Delta$, hence $M_{j \to j'}^{\mathrm{step}} = 0$.
\end{proof}

\begin{remark}[Tighter bound via mass concentration]
\label{rem:tighter-bound}
The bound $M_{j \to j'}^{\mathrm{step}} \leq \Bclip$ is loose: it
ignores that admissible $\Delta$ has total $\ell_1$ mass at most 1,
of which at most a fraction $Q^*$ lies on $A_{j, j'}$.  A tighter
bound is
\begin{equation}
\label{eq:tighter-bound}
  M_{j \to j'}^{\mathrm{step}}(\mathcal{D}) \;\leq\;
  \Bclip \cdot Q^*(\mathcal{D}),
\end{equation}
where $Q^*(\mathcal{D}) := \max_{j \neq j'} Q^*_{j, j'}(\mathcal{D})$
and
\[
  Q^*_{j, j'}(\mathcal{D}) := \sup_{\Delta \in \mathcal{D}_j,
  \|\Delta\|_1 \leq 1} \;\sum_{a \in A_{j, j'}} |\Delta(a)|.
\]
$Q^*_{j, j'}(\mathcal{D}) \in [0, 1]$ is the maximum fraction of a
unit-$\ell_1$ $j$-perturbation that lies on the shared-action set
$A_{j, j'}$.  For action-disjoint channels, $Q^*_{j,j'} = 0$.
\end{remark}

\section{Cumulative cascade-window bound}

\subsection{Aggregating per-step over $\taumeta$}

The operational target is bounding the cumulative coupling response
over the cascade-time window $\taumeta$.  Under (C11.CLK), the
maximum number of SPRT exposure events in $\taumeta$ is bounded by
$\Nev(\taumeta) \leq \lammax \cdot \taumeta$ where $\lammax$ is the
deployment-class action rate cap.

\begin{lemma}[Cumulative cascade-window coupling bound]
\label{lem:cumulative-bound}
Under (C5.HOEFF), the structural-incidence convention, and a
deployment-class action rate cap $\lammax$ (calibrated by
(C11)/Audit~7), the cumulative coupling magnitude over $\taumeta$
satisfies
\begin{equation}
\label{eq:cumulative-bound}
  \bar{M}^{\mathrm{cum}}(\mathcal{D}) \;\leq\; \Bclip \cdot
  Q^*(\mathcal{D}) \cdot \lammax \cdot \taumeta,
\end{equation}
where $\bar{M}^{\mathrm{cum}}$ is the maximum total $j'$-drift
response, summed over all events in $\taumeta$, per unit-$\ell_1$
$j$-perturbation maintained throughout the window.
\end{lemma}

\begin{proof}
At each event step $n \in \{1, \ldots, \Nev(\taumeta)\}$, the
per-step coupling response is bounded by
$M_{j \to j'}^{\mathrm{step}} \leq \Bclip \cdot Q^*$
(Remark~\ref{rem:tighter-bound}).  Summing over events:
\[
  \big| \sum_{n=1}^{\Nev(\taumeta)} \Delta\mu_{j'}^{(n)}[\Delta] \big|
  \leq \Nev(\taumeta) \cdot \Bclip \cdot Q^*
  \leq \lammax \cdot \taumeta \cdot \Bclip \cdot Q^*.
\]
Maximizing over $j \neq j'$ gives the deployment-coupling bound.
\end{proof}

\subsection{Connection to Audit 4 and (C7)}

Audit~4's empirical procedure measures the response of channel
$j'$ to a unit perturbation in channel $j$ over the cascade window
$\taumeta$.  By Lemma~\ref{lem:cumulative-bound}, the audit-time
$\bar{M}$ that operators must verify is bounded by
$\Bclip \cdot Q^* \cdot \lammax \cdot \taumeta$.

\textbf{(C7) follows as a corollary} when:
\begin{enumerate}
\item $\Bclip$ is a deployment-class constant (assumed in
  (C5.HOEFF)).
\item $Q^*(\mathcal{D})$ is intensive in $|P|$ (requires the
  structural condition (C5.OVL) below; not implied by (C5.MULT)
  alone).
\item $\lammax$ is intensive in $|P|$ (assumed in (C11) action
  rate cap).
\item $\taumeta$ is intensive in $|P|$ (deployment-class via
  \citep[Phase~Redundancy]{lasser2026phase}'s scaling).
\end{enumerate}

All four factors are deployment-class constants under their stated
conditions.  The right-hand side of \eqref{eq:cumulative-bound} is
therefore intensive in $|P|$, and $\bar{M}^{\mathrm{cum}}$ is
bounded by a $|P|$-independent constant: this is exactly what (C7)
asserts.

\section{The new structural condition: (C5.OVL) overlap-mass stability}

\subsection{Why (C5.MULT) is not sufficient}

(C5.MULT) fixes the channel partition's cardinality $\Kch$ and the
multinomial cardinality bound $\Kchmulti$ before deployment.  It
does \emph{not} prevent admissible adversaries from concentrating
$\ell_1$ mass on $A_{j, j'}$ at increasing rates as $|P|$ grows: a
more capable agent might, for fixed channel partition, find
adversarial strategies that place more probability mass on
shared-action events.  This would scale $Q^*$ with $|P|$ and
defeat the closed-form bound.

\subsection{The new sub-clause}

\begin{assumption}[Overlap-mass stability, (C5.OVL)]
\label{ass:c5-ovl}
The channel partition's shared-action mass is uniformly bounded
across capability scales: there exists a deployment-class constant
$Q_{\max} < 1$, independent of $|P|$, such that for every $j \neq
j'$,
\begin{equation}
  \sup_{|P|} \;\sup_{\Delta \in \mathcal{D}_j, \|\Delta\|_1 \leq 1}
  \;\sum_{a \in A_{j, j'}} |\Delta(a)| \;\leq\; Q_{\max}.
\end{equation}
\end{assumption}

\begin{remark}[(C5.OVL) is testable from verification-protocol design]
\label{rem:c5-ovl-design}
(C5.OVL) is a \emph{design condition} on the verification
protocol's event-classification policy.  It requires that, under
$|P|$-scaling, the shared-action sets $A_{j, j'}$ either:
\begin{enumerate}
\item Remain empty (the action-disjoint case, $Q_{\max} = 0$); or
\item Grow at rates that are matched by classification-policy
  expansion of pure-channel sets, so that the \emph{mass fraction}
  on shared actions stays bounded.
\end{enumerate}
The first case is achievable by design when the four channels of
Lemma~5b measure structurally distinct ledger event types.  The
second is operationally testable by auditing the
classification-policy's growth rule against the deployment's
capability scaling pattern.
\end{remark}

\section{Discussion}

\subsection{What the v2 closed form delivers}

If Lemma~\ref{lem:cumulative-bound} holds and (C5.OVL) is
adopted, (C7) graduates from a free-floating assumption to a
corollary of:
\begin{itemize}
\item (C5.HOEFF) per-step LLR clipping --- already required.
\item (C5.MULT) bounded channel multiplicity with structural
  incidence --- already required, slightly strengthened.
\item (C5.OVL) overlap-mass stability --- new.
\item (C11) action rate cap --- already required.
\item Phase Redundancy's $\taumeta$ scaling --- already inherited.
\end{itemize}

Three of these are existing conditions; the new content is
(C5.OVL), which encodes the structural property that the closed
form requires.  The shift in epistemic burden is from ``$\bar{M}$
is bounded by something'' to ``the verification protocol's
classification policy maintains bounded shared-action mass under
$|P|$-scaling.''  The latter is more structural and is
verification-protocol-design-controllable.

\subsection{What the v2 closed form does not deliver}

\paragraph{(C5.OVL) is a non-trivial structural condition.}  It is
not derived from existing conditions; it must be added.  It is
testable in principle (audit the classification policy's
$|P|$-scaling rule) but the test requires a deployment that has
explicitly characterized its scaling rule.  Deployments that scale
$|P|$ ad hoc (e.g., new capabilities mapped into channels by
unconstrained policy) cannot certify (C5.OVL) directly.

\paragraph{Pairwise-only bound.}  The cumulative bound is over
\emph{pairs} $(j, j')$.  For deployments where multiple channels
co-evolve simultaneously with shared-action coupling across more
than two channels, the pairwise bound may be loose: the actual
co-evolution correction term in Microfoundation's composition might
scale with a higher-order overlap quantity (triple intersections
$A_j \cap A_{j'} \cap A_{j''}$, etc.).  This is a refinement open
question.

\paragraph{Cumulative interpretation tied to $\taumeta$ as a hard
window.}  Lemma~\ref{lem:cumulative-bound} treats $\taumeta$ as a
hard upper bound on the cumulative window.  Under (C11.CLK), the
event-count $\Nev(\taumeta)$ is bounded with high probability $1 -
b_{\mathrm{clk}}$, so the cumulative-bound certification picks up
this failure probability.  This is consistent with Lemma~6's clock
comparability machinery; deployments where $b_{\mathrm{clk}}$ is
non-negligible carry the corresponding stack-failure probability in
the deployment claim.

\subsection{Connection to (C11) capability-scaling caveat}

(C11)'s capability-scaling caveat (a more capable agent may have
richer effective action space) is structurally parallel to the
$Q^*$-stability question: more capable agents may concentrate mass
on shared actions in subtler ways, defeating (C5.OVL).  The closed
form does not eliminate this concern --- it makes it explicit and
operational, by tying the worry to a measurable
verification-protocol property rather than to dynamics.

\section{Open questions for codex review}

\begin{enumerate}
\item \textbf{Is the L1-normalized definition of
  $M_{j \to j'}^{\mathrm{step}}$ in
  Definition~\ref{def:M-formal-v2} the right operational reading
  of Audit~4?}  v1's reviewer note suggests Audit~4 is windowed
  /cumulative.  v2 splits this into per-step (Definition~3) and
  cumulative (Lemma~2), with the cumulative version matching the
  audit's ``within $\taumeta$'' phrasing directly.  Is the
  L1-normalized per-step version a reasonable intermediate
  quantity, or should the proof be cumulative-only?

\item \textbf{Does Lemma~\ref{lem:per-step}'s tighter bound
  (Remark~\ref{rem:tighter-bound}) hold under all admissible
  $j$-perturbations?}  The bound $\Bclip \cdot Q^*$ uses that
  $Q^*$ measures the maximum fraction of $\ell_1$ mass on
  $A_{j, j'}$.  This implicitly assumes $\ell^{(j')}$ has worst-case
  alignment with $\Delta$; in practice the bound may be loose by a
  constant factor.  Is this tightness adequate?

\item \textbf{Is (C5.OVL) a reasonable structural condition to
  add?}  v2 adds (C5.OVL) as a new sub-clause of (C5.MULT).  The
  alternative is to treat $\bar{M}$ itself as the
  free-floating constant (current paper's posture).  Does
  (C5.OVL) provide a strictly better epistemic posture (testable
  property of verification-protocol design rather than empirical
  measurement of $M$)?

\item \textbf{Cumulative bound aggregation.}  Lemma~2 sums per-step
  responses over $\Nev(\taumeta)$ events.  This treats the
  $j$-perturbation as maintained throughout the window.  Does this
  match what Microfoundation's $K_{\mathrm{coev}}$ aggregation
  actually requires, or does $K_{\mathrm{coev}}$ aggregate via a
  different (e.g., supremum or integral) operation?  If the latter,
  the cumulative bound may need a different aggregation step.

\item \textbf{Pairwise-only bound vs.\ multi-channel co-evolution.}
  v2's bound is over channel pairs.  The cooperative-overlap
  regime in Lemma~5c involves multi-substrate dynamics where
  multiple channels co-evolve simultaneously.  Does the pairwise
  bound suffice for Microfoundation's composition correction term,
  or does the correction require a higher-order overlap bound?
\end{enumerate}

\section{Status}

v2 attempts to close the v1 P0/P1 gaps via:
\begin{itemize}
\item Lipschitz formulation of $M_{j \to j'}$ (P0-1 fix).
\item Cumulative bound directly with $\lammax \cdot \taumeta$
  factors (P0-2 fix).
\item Structural incidence map separate from LLR support (P1-4).
\item $Q_\Delta^*$ formulation over signed measures (P1-3, P1-5).
\item Explicit baseline compatibility (P1-7).
\item New structural condition (C5.OVL) for $Q^*$ intensivity
  (P1-6).
\item Pairwise terminology consistency (P2-8).
\end{itemize}

The v2 closed form is more structurally honest than v1: it
explicitly identifies the verification-protocol-design property
(C5.OVL) on which (C7)'s intensivity depends, rather than
asserting the intensivity directly.

\end{document}
